% Chapter 3 of Singular Learning Theory for the Rest of Us: the free energy of one-coefficient models.
\sltchapter{The Free Energy by Hand}{The whole theory in one coefficient}{$y=abx$, $y=ab^2x$}{ch:coefficient}
\newtonfoot{2/4}{4}{\tfrac14}{a^2b^4}
\begin{refsection}[chapters/03-coefficient.bib]

\begin{chaptersummary}
What does Bayesian model comparison charge a model whose parameters the data see only through their product? When the likelihood depends on the parameters through one scalar coefficient $c=g(w)$, the free energy is a Gaussian window of width about $n^{-1/2}$ integrated against the prior density of $c$. We compute it by hand for $y=abx$, $y=ab^2x$ and their relatives, and prove that a density $\abs{c-c_0}^{2\lambda-1}(\log\frac1{\abs{c-c_0}})^{m-1}$ at the truth gives $F_n-nS_n=\lambda\log n-(m-1)\log\log n+O_p(1)$, with an explicit limit law, convergence of means, and a correction of order $1/\log n$ when $m\ge2$. Exact quadrature up to $n=10^8$ confirms every term, the slow one included. The theorem needs a scalar coefficient, Gaussian noise and a fixed truth; Watanabe's theorem removes the first assumption at the price of resolution of singularities.
\end{chaptersummary}

\section{The question: what does a product cost?}\label{sec:coefficient:question}

A hidden unit with input weight $b$ and output weight $a$ contributes $a\tanh(bx)\approx abx$ near the point where it is switched off, and $a\relu(bx)=ab\relu(x)$ when $b>0$ (\cref{ch:singular}): there the data see the two weights only through their product. The smallest model with this feature is the \term{product model}\index{product model}
\[
y=abx+\varepsilon,\qquad x\sim\Normal(0,1),\quad\varepsilon\sim\Normal(0,1),\quad(a,b)\in[-1,1]^2,
\]
with the uniform prior, fitted to data that are pure noise, so that the true slope is $c_0=0$. A Bayesian compares models by their \term{marginal likelihoods}\index{marginal likelihood}, the probabilities that they assigned to the data before seeing them. The negative logarithm of the marginal likelihood is the \term{free energy}\index{free energy}\sidenote{The name comes from statistical physics: $e^{-F_n}$ is the normalizing constant, or partition function, of the posterior, a Gibbs measure with energy $nL_n$ at temperature $1$.}
\[
F_n=-\log\int_W\varphi(w)\prod_{i=1}^np(y_i\mid x_i,w)\dd w=-\log\int_W\varphi(w)e^{-nL_n(w)}\dd w,
\]
and $F_n-nS_n=-\log\int_W\varphi(w)e^{-nK_n(w)}\dd w$, where $nS_n$ is the negative log-likelihood of the truth, is what the data charge the model for not knowing the truth in advance. The prologue measured it for the product model (\cref{fig:prologue:free}): it grows like $\tfrac12\log n-\log\log n$. This chapter derives that formula with every constant, for every regression $y=g(w)x+\varepsilon$, whose likelihood depends on its parameters through the one number $g(w)$.

\begin{regularanswer}[the free energy]
If the model is regular with $d$ parameters, then $F_n=nL_n(\hat w)+\tfrac d2\log n+O_p(1)$, where $\hat w$ is the maximum-likelihood estimate, and since $nL_n(\hat w)-nS_n=O_p(1)$, also $F_n-nS_n=\tfrac d2\log n+O_p(1)$ \parencite{Schwarz1978}; \cref{thm:regular:laplace} gives the bounded term up to $o_p(1)$. For $y=abx$, with $d=2$, the data should charge $\log n$.
\end{regularanswer}

The regular answer, Schwarz's \term{Bayesian information criterion}\index{BIC}, comes from Laplace's method\index{Laplace approximation!failure of} (\cref{ch:regular}): each of the $d$ directions in which the posterior has width $n^{-1/2}$ contributes $\tfrac12\log n$. For $y=abx$ the method has nothing to expand around. The likelihood is maximized on the whole hyperbola $ab=\hat c$, where the Hessian of $nL_n$, the matrix $S_{xx}(b,a)^\top(b,a)$ with $S_{xx}=\sum_ix_i^2$, has rank one, so that Laplace's approximation, which contains half the logarithm of its determinant, returns $-\infty$. The Fisher information $\bigl(\begin{smallmatrix}b^2&ab\\ab&a^2\end{smallmatrix}\bigr)$ is singular everywhere and zero at the center of the cross $ab=0$ of true parameters (\cref{comp:singular:fisher,prop:singular:fisher}), the posterior spreads along the cross (\cref{thm:singular:scales}), and the prologue measured a slope of $0.39$ against $\log n$, not $1$.

The data see one number, the slope $c=ab$, so whatever replaces $\tfrac d2$ must depend on how the prior spreads that number. The question of this chapter is: \emph{what is $F_n$ for $y=abx$, exactly, and for which models does the same computation work?}

\section{By hand: one coefficient, one integral}\label{sec:coefficient:byhand}

Write $S_{xx}=\sum_ix_i^2$, $S_{xy}=\sum_ix_iy_i$ and $S_{yy}=\sum_iy_i^2$. With unit noise variance the energy of the product model is $nL_n(a,b)=\tfrac n2\log(2\pi)+\tfrac12\sum_i(y_i-abx_i)^2$, and the truth is $c_0=0$, so that $y_i=\varepsilon_i$.

\begin{computation}{The free energy of $y=abx$}\label{comp:coefficient:product}
\step \emph{The energy depends on $(a,b)$ only through $c=ab$.} Expanding the square,
\[
nL_n(a,b)=\tfrac n2\log(2\pi)+\tfrac12S_{yy}-cS_{xy}+\tfrac12c^2S_{xx}.
\]
At the truth $c=0$ this is $nS_n=\tfrac n2\log(2\pi)+\tfrac12S_{yy}$. Completing the square with $\hat c=S_{xy}/S_{xx}$ and $Z_n=S_{xy}/\sqrt{S_{xx}}$ gives
\[
nK_n(a,b)=\frac{S_{xx}}2\,(ab-\hat c)^2-\frac{Z_n^2}2.
\]
\step \emph{The data enter through two numbers.} Given the inputs, $S_{xy}=\sum_ix_i\varepsilon_i$ is normal with mean $0$ and variance $S_{xx}$. So $Z_n$ is standard normal, exactly and for every $n$, and independent of $S_{xx}$, which has the chi-square distribution with $n$ degrees of freedom. As a function of $c$, the likelihood is a Gaussian window of width $S_{xx}^{-1/2}\approx n^{-1/2}$ centered $Z_n$ widths away from the truth.
\step \emph{The density of $c=ab$.} For $\abs a$ and $\abs b$ independent and uniform on $[0,1]$, and $0<s\le1$,
\[
\Prob(\abs{ab}\le s)=\Prob(\abs a\le s)+\int_s^1\frac sa\dd a=s+s\log\frac1s.
\]
Differentiating, and splitting the mass evenly between the two signs, $c$ has the density $\rho(c)=\tfrac12\log\frac1{\abs c}$ on $[-1,1]$.
\step \emph{One integral.} Integrating $e^{-nK_n}$ against the prior is integrating the window against $\rho$:
\[
F_n-nS_n=-\frac{Z_n^2}2-\log\int_{-1}^1e^{-S_{xx}(c-\hat c)^2/2}\,\tfrac12\log\tfrac1{\abs c}\dd c.
\]
Substitute $c=t/\sqrt{S_{xx}}$, so that $\sqrt{S_{xx}}(c-\hat c)=t-Z_n$, $\dd c=\dd t/\sqrt{S_{xx}}$ and $\log\frac1{\abs c}=\tfrac12\log S_{xx}-\log\abs t$:
\begin{align*}
F_n-nS_n&=\tfrac12\log S_{xx}-\frac{Z_n^2}2\\
&\quad-\log\int_{\abs t<\sqrt{S_{xx}}}e^{-(t-Z_n)^2/2}\,\tfrac12\bigl(\tfrac12\log S_{xx}-\log\abs t\bigr)\dd t.
\end{align*}
\step \emph{The integral is elementary.} Dropping the restriction $\abs t<\sqrt{S_{xx}}$ changes the integral by at most a constant times $e^{-S_{xx}/8}\log S_{xx}$ when $\abs{Z_n}\le\frac12\sqrt{S_{xx}}$. With $U\sim\Normal(0,1)$ and $h(z)=\E\log\abs{z+U}$, the two Gaussian integrals are $\sqrt{2\pi}$ and $\sqrt{2\pi}\,h(Z_n)$, so that, up to that exponentially small error,
\[
F_n-nS_n=\tfrac12\log S_{xx}-\log\bigl(\tfrac12\log S_{xx}-h(Z_n)\bigr)-\frac{Z_n^2}2+\tfrac12\log\tfrac2\pi.
\]
\step \emph{Read off the asymptotics.} Since $S_{xx}=n(1+O_p(n^{-1/2}))$ and $h(Z_n)=O_p(1)$, the second term is $-\log\log n+\log2+O_p(1/\log n)$, and
\[
F_n-nS_n=\tfrac12\log n-\log\log n+\tfrac12\log\tfrac8\pi-\frac{Z_n^2}2+O_p\Bigl(\frac1{\log n}\Bigr).
\]
\end{computation}

The exactness of step~2 is special to Gaussian noise.\sidenote{With other noise of unit variance the Gaussian model no longer contains the truth, and $nS_n$ must be replaced by $nL_n(w_0)$, the negative log-likelihood of the model at $c=0$. The other steps hold unchanged, but $Z_n$ is standard normal only in the limit, by the central limit theorem.} At $n=10^4$, with $S_{xx}=n$ and $Z_n=0$, step~5 gives $F_n-nS_n=4.6052-\log5.2404-0.2258=2.7230$ nats, because $h(0)=\E\log\abs U=-\frac12(\gamma+\log2)=-0.6352$, with $\gamma=0.5772$ Euler's constant.\sidenote{Differentiate $\E\abs U^s=2^{s/2}\Gamma(\frac{s+1}2)/\sqrt\pi$ at $s=0$, and use $\Gamma'(\frac12)/\Gamma(\frac12)=-\gamma-2\log2$.} Step~6 gives $2.8522$; the difference is the $O_p(1/\log n)$ term. The regular plane $y=ax_1+bx_2$, with the same prior and at the same point, pays $\log n+\log\frac2\pi=8.7588$ nats: the product model, with as many parameters, is cheaper by six nats, a factor of about $420$ in marginal likelihood.

\paragraph{What one had to notice.} Three things. The likelihood depends on $(a,b)$ only through $c=ab$, so the integral over the square is an integral against the law of $c$ under the prior. The data reduce to $S_{xx}$ and $Z_n$, and $Z_n$ is exactly standard normal. And the cross of true parameters has become a logarithmic singularity of $\rho$ at $c=0$, which a window of width $n^{-1/2}$ sees as the factor $\frac12\log n$: that factor is the $-\log\log n$. The same facts hold for every model $y=g(w)x+\varepsilon$ with a scalar coefficient; only the density changes (\cref{comp:coefficient:densities}, \cref{fig:coefficient:densities}).

\begin{figure}[t]
\begin{fullwidth}
\centering
\begin{tikzpicture}
\begin{groupplot}[group style={group size=2 by 1,horizontal sep=17mm},splot,width=80mm,height=47mm]
\nextgroupplot[xlabel={coefficient $c$},ylabel={density $\rho(c)$},xmin=-0.6,xmax=0.6,ymin=0,ymax=3.2,
  xtick={-0.5,0,0.5},ytick={0,1,2,3},title={(a) densities of the coefficient},legend pos=north west,
  legend style={font=\scriptsize}]
\addplot[band,forget plot] table[x=c,y expr=2.9*\thisrow{window}] {figures/data/03-coefficient-density.dat} \closedcycle;
\addplot[regularpred] table[x=c,y=line] {figures/data/03-coefficient-density.dat};
\addplot[theory] table[x=c,y=prod] {figures/data/03-coefficient-density.dat};
\addplot[singulardeep,line width=0.9pt,densely dashed,no marks] table[x=c,y=cusp] {figures/data/03-coefficient-density.dat};
\legend{$y=ax$,$y=abx$,$y=ab^2x$}
\node[font=\scriptsize,text=singulardeep,anchor=west] at (axis cs:0.09,2.75) {window, $n=400$};
\nextgroupplot[xmode=log,ymode=log,xlabel={half-width $r$},ylabel={prior mass $\mu\{\abs c<r\}$},
  xmin=1e-8,xmax=1,ymin=1e-8,ymax=2,xtick={1e-8,1e-6,1e-4,1e-2,1},ytick={1e-8,1e-6,1e-4,1e-2,1},
  title={(b) the prior mass of a window}]
\addplot[guide,forget plot] coordinates {(1e-2,1e-8) (1e-2,2)};
\addplot[regularpred] table[x=r,y=line] {figures/data/03-coefficient-mass.dat};
\addplot[theory] table[x=r,y=prod] {figures/data/03-coefficient-mass.dat};
\addplot[theory2] table[x=r,y=chain] {figures/data/03-coefficient-mass.dat};
\addplot[singulardeep,line width=0.9pt,densely dashed,no marks] table[x=r,y=cusp] {figures/data/03-coefficient-mass.dat};
\addplot[muted,line width=0.8pt,dashdotted,no marks] table[x=r,y=cubic] {figures/data/03-coefficient-mass.dat};
\node[font=\scriptsize,text=regular,anchor=north west] at (axis cs:2e-8,1.6e-8) {$y=ax$};
\node[font=\scriptsize,text=singular,anchor=south west] at (axis cs:2e-8,4e-7) {$y=abx$};
\node[font=\scriptsize,text=singular,anchor=south west] at (axis cs:2e-8,4e-6) {$y=a_1a_2a_3x$};
\node[font=\scriptsize,text=singulardeep,anchor=south west] at (axis cs:2e-8,3e-4) {$y=ab^2x$};
\node[font=\scriptsize,text=muted,anchor=south west] at (axis cs:2e-8,4.5e-3) {$y=ab^3x$};
\node[font=\scriptsize,text=muted,anchor=south east,rotate=90] at (axis cs:1.6e-2,1e-7) {$n=10^4$};
\end{groupplot}
\end{tikzpicture}
\end{fullwidth}
\caption{The prior seen through the coefficient, for the uniform prior on $[-1,1]^d$. (a) The density $\rho$ of $c=g(w)$ for the regular line $y=ax$ (flat), the product model $y=abx$ (a logarithmic spike) and the cusp model $y=ab^2x$ (a power spike), with a likelihood window $e^{-n(c-\hat c)^2/2}$ for $n=400$ and $\hat c=0.05$ (shaded, arbitrary vertical scale). (b) The prior mass $\mu\{\abs c<r\}$ of a window of half-width $r$, from the closed forms of \cref{comp:coefficient:densities}; its slope tends to $2\lambda$: $1$ for the first three models (up to logarithms), $\frac12$ and $\frac13$ for the cusps. The free energy is about minus the logarithm of the mass at $r=n^{-1/2}$ (dashed: $n=10^4$). A Monte Carlo check with $10^7$ prior draws (seed $31$) agrees with all twenty values tested, at $r=10^{-1},\dots,10^{-4}$, within its $95\%$ Clopper--Pearson intervals. Script: \texttt{figures/scripts/03-coefficient-mass.py}.}\label{fig:coefficient:densities}
\end{figure}

\begin{computation}{Densities of coefficients}\label{comp:coefficient:densities}
\step \emph{The chain $c=a_1a_2\cdots a_L$.} For $a_i$ uniform on $[-1,1]$, $\Prob(-\log\abs{a_i}>u)=\Prob(\abs{a_i}<e^{-u})=e^{-u}$, so $-\log\abs{a_i}$ is exponential with mean $1$, and $-\log\abs c$, a sum of $L$ independent such variables, has the Gamma density $u^{L-1}e^{-u}/(L-1)!$. Changing variables to $s=e^{-u}$, and splitting the mass between the signs,
\[
\rho_L(c)=\frac1{2(L-1)!}\Bigl(\log\frac1{\abs c}\Bigr)^{L-1}\qquad(\abs c\le1).
\]
For $L=1$ this is the uniform density $\frac12$ of the regular line; for $L=2$ it is \cref{comp:coefficient:product}.
\step \emph{The cusp $c=ab^2$.} For $0<s\le1$, $\Prob(\abs ab^2\le s)=\Prob(b^2\le s)+\int_{\sqrt s}^1\frac s{b^2}\dd b=2\sqrt s-s$, so $\rho(c)=\frac12(\abs c^{-1/2}-1)$.
\step \emph{The family $c=ab^k$, $k\ge2$.} In the same way $\Prob(\abs a\abs b^k\le s)=s^{1/k}+s\int_{s^{1/k}}^1b^{-k}\dd b=\frac{ks^{1/k}-s}{k-1}$, so $\rho(c)=\frac{\abs c^{1/k-1}-1}{2(k-1)}$; for the cubic cusp $y=ab^3x$, $\rho(c)=\frac14(\abs c^{-2/3}-1)$.
\step \emph{Read off.} In each case $\rho(c)=\abs c^{2\lambda-1}(\log\frac1{\abs c})^{m-1}\ell(c)$ near $c=0$, with $\ell(c)\to\ell_0>0$ and $(\lambda,m,\ell_0)$ equal to $(\frac12,L,\frac1{2(L-1)!})$ for the chain, $(\frac14,1,\frac12)$ for the cusp, and $(\frac1{2k},1,\frac1{2(k-1)})$ for $ab^k$.
\end{computation}

\paragraph{The cusp by hand.} With the cusp's density, step~4 of \cref{comp:coefficient:product} becomes, after the same substitution and up to an exponentially small error,
\[
F_n-nS_n=\tfrac14\log S_{xx}-\frac{Z_n^2}2-\log\Bigl(\tfrac12\int_\R e^{-(t-Z_n)^2/2}\abs t^{-1/2}\dd t-\tfrac12\sqrt{2\pi}\,S_{xx}^{-1/4}\Bigr).
\]
The subtracted term is $O(n^{-1/4})$, so $F_n-nS_n=\frac14\log n+O_p(1)$, with no $\log\log n$ because the density has no logarithm. At $Z_n=0$ the integral is $2^{1/4}\Gamma(\frac14)=4.3116$, and at $n=10^4$ the free energy is $1.53$ nats ($1.59$ with the $O(n^{-1/4})$ term): cheaper still.

\paragraph{The rule, and why it holds.} In both computations the exponent of $\abs{c-c_0}$ in the density gave the coefficient of $\log n$, and the power of the logarithm gave that of $\log\log n$: a density $\abs{c-c_0}^{2\lambda-1}(\log\frac1{\abs{c-c_0}})^{m-1}$ at the truth gives $F_n-nS_n\approx\lambda\log n-(m-1)\log\log n$. Heuristically, the window keeps the coefficients within about $n^{-1/2}$ of the truth and discards the rest, so $e^{-(F_n-nS_n)}$ is, up to a bounded factor, the prior mass of the window (\cref{fig:coefficient:densities}(b)):
\[
\mu\{\abs{c-c_0}<n^{-1/2}\}\asymp\int_0^{n^{-1/2}}\delta^{2\lambda-1}\Bigl(\log\frac1\delta\Bigr)^{m-1}\dd\delta\asymp n^{-\lambda}\Bigl(\frac12\log n\Bigr)^{m-1}.
\]
Its negative logarithm is $\lambda\log n-(m-1)\log\log n+O(1)$, and \cref{sec:coefficient:theorem} turns the heuristic into a proof. This is the singular version of the Bayesian \term{Occam factor}\index{Occam factor} of \textcite[ch.~28]{MacKay2003}: not a ratio of a posterior width to a prior width, but the prior mass of a window around a singular point.\sidenote{For the product model the window keeps the region $\abs{ab}\lesssim n^{-1/2}$ around the cross, of area about $n^{-1/2}\log n$: the arms are $n^{-1/2}$ thick far from the center and fatter near it, which is where the logarithm comes from.} The exponent $\lambda$ is the \term{learning coefficient}\index{learning coefficient} and $m$ its \term{multiplicity}\index{multiplicity}; as numbers attached to $K$ and the prior they are the \term{real log canonical threshold}\index{real log canonical threshold} and its multiplicity, the same two numbers under another name (\cref{prop:coefficient:zeta} connects the definitions).

\section{The theorem: from the density to the free energy}\label{sec:coefficient:theorem}

The computations used three facts: the likelihood factors through a scalar coefficient, the data enter through a window whose position is standard normal, and the prior density of the coefficient has a power-and-logarithm singularity at the truth. We now prove the rule under exactly these hypotheses, for a \term{one-coefficient model}\index{one-coefficient model}.

\begin{assumption}[A one-coefficient regression model]\label[assumption]{ass:coefficient:onecoef}
\leavevmode
\begin{enumerate}
\item The inputs $x_1,x_2,\dots$ are independent real random variables with a common distribution, and $\E[x^2]=1$.
\item The outputs are $y_i=c_0x_i+\varepsilon_i$ for a fixed $c_0\in\R$, with $\varepsilon_1,\varepsilon_2,\dots$ independent $\Normal(0,1)$ variables, independent of the inputs.
\item The model is $p(y\mid x,w)=(2\pi)^{-1/2}\exp(-\frac12(y-g(w)x)^2)$ with a measurable $g\colon W\to\R$ and a prior probability density $\varphi$ on $W\subseteq\R^d$. The law of the coefficient $c=g(w)$ when $w\sim\varphi$ is denoted $\mu$.
\item There are $\eta\in(0,1)$, $\lambda>0$, an integer $m\ge1$, and a bounded function $\ell$ on $0<\abs{c-c_0}<\eta$ with one-sided limits $\ell_+=\lim_{c\downarrow c_0}\ell(c)$ and $\ell_-=\lim_{c\uparrow c_0}\ell(c)$ and $\ell_++\ell_->0$, such that on the interval $\abs{c-c_0}<\eta$ the law $\mu$ has the density
\[
\rho(c)=\abs{c-c_0}^{2\lambda-1}\Bigl(\log\frac1{\abs{c-c_0}}\Bigr)^{m-1}\ell(c)\qquad(0<\abs{c-c_0}<\eta).
\]
\end{enumerate}
\end{assumption}

The two-sided case is $\ell$ continuous and positive at $c_0$; the case $\ell_-=0$ covers a coefficient that cannot fall below the truth, such as $c=a^2$ at $c_0=0$ (\cref{ex:coefficient:onesided}). Nothing is assumed about $\mu$ away from $c_0$, and neither compactness of $W$ nor analyticity of $g$, which Watanabe's general theorem assumes, is needed.\sidenote{A prior with an atom at $c_0$, such as a spike-and-slab prior that switches the coefficient off with probability $p>0$, violates part~4. Then the integral in \eqref{eq:coefficient:exact} below is at least $pe^{-Z_n^2/2}$, and $F_n-nS_n\le\log\frac1p$ stays bounded.} The inputs need only a second moment; the noise must be Gaussian of known variance, which makes the energy exactly quadratic in $c$. Write
\[
S_{xx}=\sum_{i=1}^nx_i^2,\qquad T_n=\sum_{i=1}^nx_i\varepsilon_i,\qquad\hat c=c_0+\frac{T_n}{S_{xx}},\qquad Z_n=\frac{T_n}{\sqrt{S_{xx}}}
\]
on the event $\{S_{xx}>0\}$, and $Z_n=0$ on $\{S_{xx}=0\}$; $\hat c=S_{xy}/S_{xx}$ is the least-squares slope.

\begin{lemma}[Exact representation]\label[lemma]{lem:coefficient:exact}
Assume parts 1--3 of \cref{ass:coefficient:onecoef}. On the event $\{S_{xx}>0\}$, for every $w\in W$,
\begin{equation}\label{eq:coefficient:energy}
nK_n(w)=\frac{S_{xx}}2\bigl(g(w)-\hat c\bigr)^2-\frac{Z_n^2}2,
\end{equation}
and
\begin{equation}\label{eq:coefficient:exact}
F_n-nS_n=-\frac{Z_n^2}2-\log\int_\R e^{-S_{xx}(c-\hat c)^2/2}\,\mu(\dd c).
\end{equation}
On $\{S_{xx}=0\}$, $F_n-nS_n=0$. Conditionally on $x_1,\dots,x_n$, the variable $Z_n$ is standard normal on $\{S_{xx}>0\}$; if $\Prob(x=0)=0$, then $Z_n\sim\Normal(0,1)$ is independent of $S_{xx}$.
\end{lemma}

\begin{proof}
The plan is to expand the squares, complete the square in the coefficient, and push the prior forward to the coefficient. Write $c=g(w)$. Since $y_i-cx_i=\varepsilon_i-(c-c_0)x_i$, expanding the squares in $nL_n(w)=\frac n2\log(2\pi)+\frac12\sum_i(y_i-cx_i)^2$ and in $nS_n=\frac n2\log(2\pi)+\frac12\sum_i\varepsilon_i^2$ gives $nK_n(w)=\frac12S_{xx}(c-c_0)^2-(c-c_0)T_n$. On $\{S_{xx}>0\}$, completing the square gives \eqref{eq:coefficient:energy}, because $\hat c-c_0=T_n/S_{xx}$ and $T_n^2/S_{xx}=Z_n^2$; on $\{S_{xx}=0\}$ every $x_i$ vanishes, so $K_n\equiv0$ and $F_n-nS_n=0$. The integrand of $F_n-nS_n=-\log\int_W\varphi e^{-nK_n}$ depends on $w$ only through $g(w)$, and $\int_Wf(g(w))\varphi(w)\dd w=\int_\R f(c)\,\mu(\dd c)$ for every measurable $f\ge0$, by the definition of $\mu$; with $f(c)=\exp(-S_{xx}(c-\hat c)^2/2+Z_n^2/2)$ this is \eqref{eq:coefficient:exact}. Finally, given the inputs, $T_n$ is a fixed linear combination of the independent $\Normal(0,1)$ variables $\varepsilon_i$, so it is $\Normal(0,S_{xx})$, and $Z_n$ is $\Normal(0,1)$ on $\{S_{xx}>0\}$. If $\Prob(x=0)=0$, then $S_{xx}>0$ almost surely, and since this conditional law does not depend on the inputs, $Z_n$ is independent of them.
\end{proof}

The lemma reduces the free energy to one dimension, at every $n$. The next lemma is the analysis: what a \term{Gaussian window}\index{Gaussian window} of width $s^{-1/2}$ sees when it is placed $z$ widths from the singular point of $\rho$.

\begin{lemma}[A Gaussian window on a singular density]\label[lemma]{lem:coefficient:window}
Let $\rho$, $\eta$, $\lambda$, $m$, $\ell$ and $\ell_\pm$ be as in part~4 of \cref{ass:coefficient:onecoef}, and write $\ell_{\sgn t}$ for $\ell_+$ when $t>0$ and for $\ell_-$ when $t<0$. For $s>1$ and $z\in\R$ put
\begin{align*}
J_s(z)&=s^{\lambda}\bigl(\tfrac12\log s\bigr)^{1-m}\int_{\abs{c-c_0}<\eta}e^{-s(c-c_0-z/\sqrt s)^2/2}\,\rho(c)\dd c,\\
J(z)&=\int_\R e^{-(t-z)^2/2}\,\abs t^{2\lambda-1}\ell_{\sgn t}\dd t.
\end{align*}
Then $J$ is continuous and positive, and $\sup_{\abs z\le M}\abs{J_s(z)-J(z)}\to0$ as $s\to\infty$, for every $M>0$. Moreover, there is a constant $C$ such that $J_s(z)\le C(1+\abs z)^{(2\lambda-1)_+}$ for all $s\ge e^2$ and $z\in\R$.
\end{lemma}

\begin{proof}
The plan is to rescale the coefficient by $\sqrt s$, which turns the window into a fixed Gaussian and the density into a function that converges pointwise, and then to use dominated convergence.

\emph{Rescaling.} Substitute $c=c_0+t/\sqrt s$. Then $\abs{c-c_0}^{2\lambda-1}=s^{1/2-\lambda}\abs t^{2\lambda-1}$, $\dd c=s^{-1/2}\dd t$, and $\log\frac1{\abs{c-c_0}}=\frac12\log s\,(1+u_s(t))$ with $u_s(t)=\log\frac1{\abs t}\big/\frac12\log s$. The powers of $s$ cancel, and $J_s(z)=\int_\R e^{-(t-z)^2/2}\abs t^{2\lambda-1}f_s(t)\dd t$ with
\[
f_s(t)=\ind\{0<\abs t<\eta\sqrt s\}\,\bigl(1+u_s(t)\bigr)^{m-1}\,\ell\bigl(c_0+t/\sqrt s\bigr).
\]

\emph{A bound on $f_s$.} Let $\bar\ell=\sup\ell$, and let $s\ge e^2$, so that $\frac12\log s\ge1$. If $0<\abs t<\eta\sqrt s$, then $\log\frac1{\abs t}>-\frac12\log s$, because $\eta<1$, so $1+u_s(t)>0$; moreover $1+u_s(t)\le1$ when $\abs t\ge1$, and $1+u_s(t)\le1+\log\frac1{\abs t}$ when $\abs t<1$. With $\log_+v=\max(\log v,0)$, this gives
\[
0\le f_s(t)\le\bar\ell\,b(t),\qquad b(t)=\bigl(1+\log_+\tfrac1{\abs t}\bigr)^{m-1}.
\]

\emph{Convergence.} Fix $t\ne0$. As $s\to\infty$, eventually $\abs t<\eta\sqrt s$, $u_s(t)\to0$, and $c_0+t/\sqrt s$ tends to $c_0$ from the side of $\sgn t$; so $f_s(t)\to\ell_{\sgn t}$. For $\abs z\le M$ and every $t$, $(t-z)^2\ge(\abs t-M)_+^2$, so
\[
\sup_{\abs z\le M}\abs{J_s(z)-J(z)}\le\int_\R e^{-(\abs t-M)_+^2/2}\abs t^{2\lambda-1}\abs{f_s(t)-\ell_{\sgn t}}\dd t.
\]
The integrand tends to $0$ for every $t\ne0$. It is at most $2\bar\ell$ times the function $e^{-(\abs t-M)_+^2/2}\abs t^{2\lambda-1}b(t)$, which is integrable because $2\lambda-1>-1$, because $b$ grows only like a power of $\log\frac1{\abs t}$ near $0$, and because of the Gaussian factor. Dominated convergence gives the uniform convergence; the same domination shows that $J$ is continuous, and $J>0$ because $\ell_\pm\ge0$ and $\ell_++\ell_->0$.

\emph{The bound.} By the bound on $f_s$, $J_s(z)\le\bar\ell\int_\R e^{-(t-z)^2/2}\abs t^{2\lambda-1}b(t)\dd t$. The part over $\abs t\le1$ is at most $\int_{\abs t\le1}\abs t^{2\lambda-1}b(t)\dd t<\infty$. On $\abs t>1$, $b=1$; if $2\lambda-1\le0$ that part is at most $\sqrt{2\pi}$, and if $p=2\lambda-1>0$, then $\abs t^p\le2^p(\abs{t-z}^p+\abs z^p)$ makes it at most $2^p\bigl(\int e^{-v^2/2}\abs v^p\dd v+\sqrt{2\pi}\abs z^p\bigr)$. Both cases give the bound $C(1+\abs z)^{(2\lambda-1)_+}$.
\end{proof}

\begin{theorem}[The free energy of a one-coefficient model]\label{thm:coefficient:free}
Under \cref{ass:coefficient:onecoef}, let $R_n=F_n-nS_n-\lambda\log n+(m-1)\log\log n$, and for $z\in\R$ let
\begin{equation}\label{eq:coefficient:chi}
\chi(z)=(m-1)\log2-\log\int_\R e^{-t^2/2+zt}\,\abs t^{2\lambda-1}\ell_{\sgn t}\dd t.
\end{equation}
Let $Z_n$ be as in \cref{lem:coefficient:exact}, which is standard normal given the inputs on an event whose probability tends to one. Then:
\begin{enumerate}
\item $R_n-\chi(Z_n)\to0$ in probability. Consequently
\[
F_n-nS_n=\lambda\log n-(m-1)\log\log n+O_p(1),
\]
and $R_n$ converges in law to $\chi(Z)$, where $Z\sim\Normal(0,1)$.
\item $\E[R_n]\to\E[\chi(Z)]$, and $\E\abs{\chi(Z)}<\infty$.
\end{enumerate}
\end{theorem}

\begin{proof}
The plan is to split the integral in \eqref{eq:coefficient:exact} into the window $\abs{c-c_0}<\eta$, which \cref{lem:coefficient:window} controls, and the rest, which the Gaussian factor makes negligible; part~2 then follows from bounds that make $(R_n)$ uniformly integrable. Write $S=S_{xx}$.

\emph{Step 1: decomposition.} On $\{S>1\}$ write the integral in \eqref{eq:coefficient:exact} as $A_n+B_n$, with $A_n$ the part over $\abs{c-c_0}<\eta$. Since $\hat c=c_0+Z_n/\sqrt S$, the definition of $J_s$ gives $A_n=S^{-\lambda}(\frac12\log S)^{m-1}J_S(Z_n)$, and the identity $e^{-t^2/2+zt}=e^{z^2/2}e^{-(t-z)^2/2}$ turns \eqref{eq:coefficient:chi} into $\chi(z)=(m-1)\log2-\frac{z^2}2-\log J(z)$. Inserting $A_n$ into \eqref{eq:coefficient:exact} and using $\log(\frac12\log S)-\log\log n=\log\frac{\log S}{\log n}-\log2$, we get on $\{S>1\}$
\begin{equation}\label{eq:coefficient:decomp}
R_n=\chi(Z_n)+\lambda\log\frac Sn-(m-1)\log\frac{\log S}{\log n}-\log\frac{J_S(Z_n)}{J(Z_n)}-\log\Bigl(1+\frac{B_n}{A_n}\Bigr).
\end{equation}

\emph{Step 2: the corrections vanish.} By the strong law of large numbers and part~1 of \cref{ass:coefficient:onecoef}, $S/n\to1$ almost surely, so $\Prob(S>1)\to1$ and the second and third terms of \eqref{eq:coefficient:decomp} tend to $0$. Fix $\epsilon>0$ and choose $M$ with $\Prob(\abs{Z_n}>M)\le\epsilon$ for all $n$, which is possible because $Z_n$ is standard normal or $0$ given the inputs. With $j_M=\min_{\abs z\le M}J(z)>0$, on $\{\abs{Z_n}\le M\}$ we have $\abs{J_S(Z_n)/J(Z_n)-1}\le j_M^{-1}\sup_{\abs z\le M}\abs{J_S(z)-J(z)}$, which tends to $0$ almost surely by \cref{lem:coefficient:window}. For the last term, if $\abs{\hat c-c_0}\le\eta/2$ then $\abs{c-\hat c}\ge\eta/2$ whenever $\abs{c-c_0}\ge\eta$, so $B_n\le e^{-S\eta^2/8}$, because $\mu$ is a probability measure; and $A_n\ge S^{-\lambda}J_S(Z_n)$ when $S\ge e^2$. For large $n$ the event $\{n/2\le S\le2n,\ \abs{Z_n}\le M,\ J_S(Z_n)\ge j_M/2\}$ has probability at least $1-2\epsilon$, and on it $\abs{\hat c-c_0}=\abs{Z_n}/\sqrt S\le\eta/2$ and $B_n/A_n\le2j_M^{-1}(2n)^\lambda e^{-n\eta^2/16}\to0$. As $\epsilon$ was arbitrary, $R_n-\chi(Z_n)\to0$ in probability. Since $Z_n$ is standard normal given the inputs on $\{S>0\}$ and $\Prob(S>0)\to1$, $\chi(Z_n)$ converges in law to $\chi(Z)$, and so does $R_n$, by Slutsky's lemma (\cref{def:toolkit:convergence}); a sequence that converges in law is bounded in probability. This proves part~1.

\emph{Step 3: an upper bound.} Part~2 follows once $(R_n)$ is uniformly integrable. Suppose $\ell_+>0$ (the case $\ell_->0$ is the mirror image), and shrink $\eta$ so that $\ell\ge\ell_+/2$ on $(c_0,c_0+\eta)$. Let $n>\eta^{-2}$ and $r=n^{-1/2}$. For $c_0<c<c_0+r$, the inequality $(c-\hat c)^2\le2(c-c_0)^2+2(\hat c-c_0)^2$ gives $S(c-\hat c)^2/2\le S/n+Z_n^2$, and since $\log\frac1\delta\ge\log\frac1r$ for $\delta\le r$,
\[
\mu\{c_0<c<c_0+r\}\ge\frac{\ell_+}2\Bigl(\log\frac1r\Bigr)^{m-1}\int_0^r\delta^{2\lambda-1}\dd\delta=\frac{\ell_+}{4\lambda}\,n^{-\lambda}\bigl(\tfrac12\log n\bigr)^{m-1}.
\]
So the integral in \eqref{eq:coefficient:exact} is at least $e^{-S/n-Z_n^2}$ times this mass, and on $\{S>0\}$
\begin{equation}\label{eq:coefficient:upper}
R_n\le\frac{Z_n^2}2+\frac Sn+C_1,\qquad C_1=(m-1)\log2+\log\frac{4\lambda}{\ell_+};
\end{equation}
on $\{S=0\}$, $R_n=-\lambda\log n+(m-1)\log\log n$, which is negative for large $n$.

\emph{Step 4: a lower bound.} The integral in \eqref{eq:coefficient:exact} is at most $\mu(\R)=1$, so $R_n\ge-Z_n^2/2-\lambda\log n$ always. On $G_n=\{S\ge n/2,\ \abs{Z_n}\le\eta\sqrt S/2\}$ we do better: there $B_n\le e^{-n\eta^2/16}$, as in Step~2, and $A_n\le C(1+\abs{Z_n})^pS^{-\lambda}(\frac12\log S)^{m-1}$ with $p=(2\lambda-1)_+$, by \cref{lem:coefficient:window}. Since $s\mapsto s^{-\lambda}(\log s)^{m-1}$ decreases for $s\ge e^{(m-1)/\lambda}$, for large $n$ and $S\ge n/2$ we have $S^{-\lambda}(\frac12\log S)^{m-1}\le2^\lambda n^{-\lambda}(\frac12\log n)^{m-1}$, and also $e^{-n\eta^2/16}\le n^{-\lambda}(\frac12\log n)^{m-1}$. Hence, for large $n$, on $G_n$,
\begin{equation}\label{eq:coefficient:lower}
R_n\ge-\frac{Z_n^2}2-p\log(1+\abs{Z_n})-C_2,\qquad C_2=\log(2^\lambda C+1)-(m-1)\log2.
\end{equation}

\emph{Step 5: uniform integrability.} By \eqref{eq:coefficient:upper}, \eqref{eq:coefficient:lower} and Step~4, for large $n$,
\[
\abs{R_n}\le Y_n+\lambda\log n\;\ind_{G_n^c},\qquad Y_n=\frac{Z_n^2}2+\frac Sn+p\log(1+\abs{Z_n})+\abs{C_1}+\abs{C_2}.
\]
The family $(Y_n)$ is uniformly integrable, because $\Prob(\abs{Z_n}>u)\le\Prob(\abs Z>u)$ for every $u>0$ and $S/n\to1$ in mean by the law of large numbers in $L^1$. The second term has mean $\lambda\log n\,\Prob(G_n^c)\to0$, because $\Prob(G_n^c)$ decays exponentially: given the inputs on $\{S\ge n/2\}$, $\Prob(\abs{Z_n}>\eta\sqrt S/2)\le2e^{-\eta^2n/16}$, and $\Prob(S<n/2)\le e^{\theta n/2}(\E e^{-\theta x^2})^n\le e^{-\theta n/4}$ for a small $\theta>0$ with $\E[e^{-\theta x^2}]\le1-\frac34\theta$, which exists because $(1-e^{-\theta x^2})/\theta$ increases to $x^2$ as $\theta\downarrow0$ (monotone convergence). Nonnegative variables whose means tend to $0$ form a uniformly integrable sequence, and so does a sum of two such sequences. Hence $(R_n)$ is uniformly integrable, convergence in law gives $\E[R_n]\to\E[\chi(Z)]$, and $\E\abs{\chi(Z)}\le\sup_n\E\abs{R_n}<\infty$.
\end{proof}

The limit laws of the running models are values of Kummer's function\index{Kummer's function} ${}_1F_1(a;b;v)=\sum_{k\ge0}\frac{(a)_k}{(b)_k}\frac{v^k}{k!}$, where $(a)_k=a(a+1)\cdots(a+k-1)$.

\begin{corollary}[The running models]\label[corollary]{cor:coefficient:running}
Let the inputs and the noise be as in parts 1--2 of \cref{ass:coefficient:onecoef}, let the prior be uniform on $[-1,1]^d$, and let $c_0=0$. For the five models of \cref{tab:coefficient:running}, \cref{ass:coefficient:onecoef} holds with the listed $\lambda$, $m$ and $\ell_+=\ell_-$, and the limit law of \cref{thm:coefficient:free} is
\[
\chi(z)=(m-1)\log2-\log\Bigl(\ell_+\,2^\lambda\,\Gamma(\lambda)\;{}_1F_1\bigl(\lambda;\tfrac12;\tfrac{z^2}2\bigr)\Bigr).
\]
For $\lambda=\frac12$ this is $(m-1)\log2-\log(\ell_+\sqrt{2\pi})-\frac{z^2}2$. In particular, for the product model,
\begin{gather*}
F_n-nS_n=\tfrac12\log n-\log\log n+\tfrac12\log\tfrac8\pi-\tfrac12Z_n^2+o_p(1),\\
\E[F_n-nS_n]=\tfrac12\log n-\log\log n+\tfrac12\log\tfrac8\pi-\tfrac12+o(1).
\end{gather*}
\end{corollary}

\begin{table}[t]
\begin{fullwidth}
\centering\small
\setlength{\tabcolsep}{5pt}
\begin{tabular}{@{}llllcccS[table-format=-1.3]S[table-format=-1.3]c@{}}
\toprule
model & coefficient & $\rho(c)$, $0<\abs c<1$ & $\ell(c)$ & $\lambda$ & $m$ & $\ell_\pm$ & {$\chi(0)$} & {$\E\chi(Z)$} & $\tfrac d2$\\
\midrule
$y=ax$ & $a$ & $\frac12$ & $\frac12$ & $\frac12$ & $1$ & $\frac12$ & -0.226 & -0.726 & $\frac12$\\
$y=abx$ & $ab$ & $\frac12\log\frac1{\abs c}$ & $\frac12$ & $\frac12$ & $2$ & $\frac12$ & 0.467 & -0.033 & $1$\\
$y=a_1a_2a_3x$ & $a_1a_2a_3$ & $\frac14\log^2\frac1{\abs c}$ & $\frac14$ & $\frac12$ & $3$ & $\frac14$ & 1.854 & 1.354 & $\frac32$\\
$y=ab^2x$ & $ab^2$ & $\frac12(\abs c^{-1/2}-1)$ & $\frac12(1-\abs c^{1/2})$ & $\frac14$ & $1$ & $\frac12$ & -0.768 & -1.073 & $1$\\
$y=ab^3x$ & $ab^3$ & $\frac14(\abs c^{-2/3}-1)$ & $\frac14(1-\abs c^{2/3})$ & $\frac16$ & $1$ & $\frac14$ & -0.446 & -0.671 & $1$\\
\bottomrule
\end{tabular}
\end{fullwidth}
\caption{The running one-coefficient models at the truth $0$, with the uniform prior on $[-1,1]^d$: the density of the coefficient, the factor $\ell$ of \cref{ass:coefficient:onecoef}, the learning coefficient and multiplicity, and the limit law of $R_n$ (\cref{cor:coefficient:running}). The value $\chi(0)$ is also the constant of the population free energy $-\log\int e^{-nK}\varphi=\lambda\log n-(m-1)\log\log n+\chi(0)+o(1)$. The means $\E\chi(Z)$ of the cusps are computed by $80$-node Gauss--Hermite quadrature in \texttt{03-coefficient-free.py}; the others are $\chi(0)-\frac12$.}\label{tab:coefficient:running}
\end{table}

\begin{proof}
The plan is to read the hypotheses off the densities computed by hand and to evaluate \eqref{eq:coefficient:chi} as a power series. By \cref{comp:coefficient:product,comp:coefficient:densities}, in each case $\ell$ is bounded on $0<\abs c<1$ and tends to the listed value as $c\to0$, so \cref{ass:coefficient:onecoef} holds with every $\eta<1$. For $\ell_+=\ell_-$ the integral in \eqref{eq:coefficient:chi} is $2\ell_+\int_0^\infty t^{2\lambda-1}e^{-t^2/2}\cosh(zt)\dd t$. Expanding $\cosh(zt)$, integrating term by term (monotone convergence), and using $\int_0^\infty t^{a-1}e^{-t^2/2}\dd t=2^{a/2-1}\Gamma(\frac a2)$ (\cref{def:toolkit:gamma}, after $\tau=t^2/2$), $\Gamma(\lambda+k)=\Gamma(\lambda)(\lambda)_k$ and $(2k)!=4^kk!(\frac12)_k$ gives
\[
\ell_+2^\lambda\sum_{k\ge0}\Gamma(\lambda+k)\frac{2^kz^{2k}}{(2k)!}=\ell_+2^\lambda\Gamma(\lambda)\sum_{k\ge0}\frac{(\lambda)_k}{(\frac12)_k}\frac{(z^2/2)^k}{k!}.
\]
For $\lambda=\frac12$ the series is $e^{z^2/2}$; for the product model $(m-1)\log2-\log(\frac12\sqrt{2\pi})=\frac12\log\frac8\pi$, and part~2 of \cref{thm:coefficient:free} with $\E[Z^2/2]=\frac12$ gives the mean.
\end{proof}

The bronze curve of the prologue's \cref{fig:prologue:free} is this corollary at finite $n$: the formula of step~4 of \cref{comp:coefficient:product} with $S_{xx}$ replaced by $n$, averaged over $Z_n$; its limit $\frac12\log n-\log\log n-\frac12+\log\frac4{\sqrt{2\pi}}$ is the corollary's, since $\log\frac4{\sqrt{2\pi}}=\frac12\log\frac8\pi$. For the chain of every length $L$ the same proof gives $\chi(z)=\frac12\log\frac{2^{2L-1}((L-1)!)^2}\pi-\frac{z^2}2$.

\section{The theorem: the limit law and its slow approach}\label{sec:coefficient:limit}

The limit law is a function of the position of the window in units of its width; its $-\frac{z^2}2$ part is the energy $nK_n=-Z_n^2/2$ of the least-squares fit. It is a limit in law only: $R_n$ does not converge in probability, because $Z_n$ and $Z_{4n}$ are correlated but not equal (\cref{ex:coefficient:notinprob}). And when $m\ge2$ it is approached slowly, because the logarithm in $\rho$ brings a correction of order $1/\log n$.

\begin{proposition}[The approach to the limit]\label[proposition]{prop:coefficient:rate}
Assume \cref{ass:coefficient:onecoef}, $\E[x^4]<\infty$, and $\abs{\ell(c)-\ell_{\sgn(c-c_0)}}\le C\abs{c-c_0}^\alpha$ for $0<\abs{c-c_0}<\eta$, with some $C,\alpha>0$. Let $\kappa(z)$ be the mean of $\log\abs t$ under the probability density on $\R$ proportional to $e^{-(t-z)^2/2}\abs t^{2\lambda-1}\ell_{\sgn t}$. Then
\begin{equation}\label{eq:coefficient:rate}
R_n=\chi(Z_n)+\frac{2(m-1)}{\log n}\,\kappa(Z_n)+O_p\Bigl(\frac1{(\log n)^2}\Bigr).
\end{equation}
If $\lambda=\frac12$ and $\ell_+=\ell_-$, then $\kappa(z)=\E\log\abs{z+U}$ with $U\sim\Normal(0,1)$, and $\E[\kappa(Z)]=-\gamma/2$ for $Z\sim\Normal(0,1)$.
\end{proposition}

\begin{proof}
The plan is to expand the factor $(1+u_s)^{m-1}$ in the proof of \cref{lem:coefficient:window} to first order in $1/\log s$, uniformly over $\abs z\le M$, and to return to \eqref{eq:coefficient:decomp}. Write $H=\frac12\log s\ge1$ and $u=\log\frac1{\abs t}$, so that $(1+u_s(t))^{m-1}=1+(m-1)u/H+r$ with $\abs r\le\sum_{k\ge2}\binom{m-1}k\abs u^kH^{-k}\le(1+\abs u)^{m-1}/H^2$, and write $\ell(c_0+t/\sqrt s)=\ell_{\sgn t}+e(t)$ with $\abs{e(t)}\le C\abs t^\alpha s^{-\alpha/2}$. With $J_1(z)=\int_\R\log\abs t\;e^{-(t-z)^2/2}\abs t^{2\lambda-1}\ell_{\sgn t}\dd t$, inserting both expansions into the formula for $J_s$ gives, uniformly over $\abs z\le M$,
\[
J_s(z)=J(z)-\frac{m-1}H\,J_1(z)+O\Bigl(\frac1{H^2}+s^{-\alpha/2}\Bigr):
\]
the error terms from $r$, from $e$, and from the region $\abs t\ge\eta\sqrt s$ that $f_s$ excludes are bounded by $H^{-2}$, by $s^{-\alpha/2}$ and by $e^{-(\eta\sqrt s-M)^2/4}$, each times an integral $\int e^{-(\abs t-M)_+^2/4}\abs t^{2\lambda-1}(1+\abs{\log\abs t})^{m-1}(1+\abs t^\alpha)\dd t<\infty$. Since $J\ge j_M>0$ and $\kappa=J_1/J$ is bounded on $\abs z\le M$, taking logarithms gives $-\log(J_s(z)/J(z))=(m-1)\kappa(z)/H+O(H^{-2}+s^{-\alpha/2})$ uniformly there.

In \eqref{eq:coefficient:decomp}, Chebyshev's inequality and $\E[x^4]<\infty$ give $S/n=1+O_p(n^{-1/2})$, so the second term is $O_p(n^{-1/2})$, the third is $O_p(n^{-1/2}/\log n)$, and $2/\log S=2/\log n+O_p(n^{-1/2}(\log n)^{-2})$; Step~2 of the proof of \cref{thm:coefficient:free} showed that the last term is smaller than every power of $n$ with probability as close to one as we like. Since $\Prob(\abs{Z_n}>M)$ is small uniformly in $n$ for large $M$, the expansion at $z=Z_n$ gives \eqref{eq:coefficient:rate}. If $\lambda=\frac12$ and $\ell_+=\ell_-$, the density defining $\kappa$ is that of $\Normal(z,1)$, so $\kappa(z)=\E\log\abs{z+U}$; for independent $Z$ and $U$, $Z+U$ has the law of $\sqrt2\,V$ with $V\sim\Normal(0,1)$, and $\E[\kappa(Z)]=\frac12\log2-\frac12(\gamma+\log2)=-\frac\gamma2$.
\end{proof}

For the product model and every chain the first-order term has mean $-(m-1)\gamma/\log n$: $-0.063$ at $n=10^4$ and $-0.031$ at $n=10^8$ for $y=abx$, twice that for the chain of length three, whose mean remainder, by the finite-$n$ prediction of \cref{sec:coefficient:measuring}, comes within $0.05$ of its limit only for $n>10^{11}$. We proved the expansion in probability only; for the means it would need the uniform integrability of $(\log n)^2(R_n-\chi(Z_n))$, which we do not prove, and \cref{sec:coefficient:measuring} checks it numerically.

\paragraph{The coefficient density is the state density.} The exponents of $\rho$ are the learning coefficient and multiplicity in the sense of Part~II, where they are defined through the volume of nearly true parameters and the poles of the \term{zeta function}\index{zeta function} $\zeta(z)=\int K(w)^z\varphi(w)\dd w$ (\cref{ch:volume}).\sidenote{Physicists would call $\rho$, or the density of $K(w)$ under the prior, a density of states; Watanabe calls the latter the state density.} For one coefficient the definitions can be compared directly.

\begin{proposition}[Density, volume and zeta function]\label[proposition]{prop:coefficient:zeta}
Assume parts 3 and 4 of \cref{ass:coefficient:onecoef}, and let $K(w)=\frac12(g(w)-c_0)^2$, the population counterpart of $K_n$ when $\E[x^2]=1$. Then:
\begin{enumerate}
\item the \term{volume function}\index{volume function} $V(\eps)=\int_{K(w)<\eps}\varphi(w)\dd w$ equals $\mu\{\abs{c-c_0}<\sqrt{2\eps}\}$, and $V(\eps)\sim\lambda^{-1}2^{\lambda-m}(\ell_++\ell_-)\,\eps^\lambda(\log\frac1\eps)^{m-1}$ as $\eps\to0$;
\item the zeta function is finite for $-\lambda<z\le0$ and infinite at $z=-\lambda$, and $(z+\lambda)^m\zeta(z)\to2^{\lambda-m}(m-1)!\,(\ell_++\ell_-)$ as $z\downarrow-\lambda$.
\end{enumerate}
\end{proposition}

\begin{proof}
The plan is to reduce both statements to one integral near $c_0$. The definition of $\mu$ gives the formula for $V$, and $\zeta(z)=2^{-z}\int_\R\abs{c-c_0}^{2z}\mu(\dd c)$. With $\tilde\ell(\delta)=\ell(c_0+\delta)+\ell(c_0-\delta)$, put $I(\sigma,r)=\int_0^r\delta^{2\sigma-1}(\log\frac1\delta)^{m-1}\tilde\ell(\delta)\dd\delta$ for $\sigma>0$ and $0<r\le\eta$. For part~1, $V(\eps)=I(\lambda,r)$ with $r=\sqrt{2\eps}$, and for $r\le e^{-1}$ the substitution $\delta=rv$ gives
\[
I(\lambda,r)=r^{2\lambda}\Bigl(\log\frac1r\Bigr)^{m-1}\int_0^1v^{2\lambda-1}\Bigl(1+\frac{\log(1/v)}{\log(1/r)}\Bigr)^{m-1}\tilde\ell(rv)\dd v.
\]
As $r\to0$ the integral tends to $(\ell_++\ell_-)/(2\lambda)$ by dominated convergence, with the bound $2\bar\ell\,v^{2\lambda-1}(1+\log\frac1v)^{m-1}$, and $r^{2\lambda}(\log\frac1r)^{m-1}\sim(2\eps)^\lambda2^{1-m}(\log\frac1\eps)^{m-1}$. For part~2, $\zeta(z)$ is $2^{-z}I(z+\lambda,\eta)$ plus a part from $\abs{c-c_0}\ge\eta$ that is at most $(\eta^2/2)^z$ for $z\le0$, and $I(\sigma,\eta)$ is finite exactly when $\sigma>0$, because $\tilde\ell\ge\frac12(\ell_++\ell_-)>0$ near $0$. The substitutions $\delta=e^{-v}$ and $y=2\sigma v$ give
\[
(2\sigma)^mI(\sigma,\eta)=\int_{2\sigma\log(1/\eta)}^\infty e^{-y}y^{m-1}\,\tilde\ell\bigl(e^{-y/(2\sigma)}\bigr)\dd y\to(m-1)!\,(\ell_++\ell_-)
\]
as $\sigma\downarrow0$, by dominated convergence. With $\sigma=z+\lambda$, multiplying by $2^{-z}\to2^\lambda$ and dividing by $2^m$ gives the limit; the part from $\abs{c-c_0}\ge\eta$ disappears after multiplication by $\sigma^m$.
\end{proof}

So $\lambda$ is the exponent of the volume and the abscissa of convergence of $\zeta$, and $m$ the power of $\log\frac1\eps$ in the volume and the order of the blow-up of $\zeta$ at $-\lambda$. For the running models $\zeta$ is meromorphic with its largest pole at $-\lambda$, of order $m$: $2^{-z}(2z+1)^{-L}$ for the chain (\cref{ex:coefficient:products}), $2^{-z}(4z+1)^{-1}(2z+1)^{-1}$ for the cusp, $2^{-z}(6z+1)^{-1}(2z+1)^{-1}$ for the cubic cusp. For analytic $K$ in general the meromorphic continuation goes back to \textcite{BernsteinGelfand1969} and \textcite{Atiyah1970}, and \cref{thm:volume:zeta} relates the volume and the zeta function. In one coefficient, $\rho$ determines Watanabe's \term{state density}\index{state density}: $v(t)=V'(t)=\bigl(\rho(c_0+\sqrt{2t})+\rho(c_0-\sqrt{2t})\bigr)/\sqrt{2t}$.

\subsection{What Watanabe's theorem adds}\label{sec:coefficient:general}

For a general model the free energy is a $d$-dimensional integral, and Watanabe's theorem says that the conclusion of \cref{thm:coefficient:free} survives: if the model is realizable and satisfies his analytic conditions (\cref{ass:toolkit:fundamental}), and $\lambda$ and $m$ are the learning coefficient and multiplicity of $K$, then $F_n=nS_n+\lambda\log n-(m-1)\log\log n+R_n$ with $R_n$ converging in law. This is Theorem~2 of \textcite{Watanabe2013}, who credits the expectation version to \textcite{Watanabe2001a} and the version as a random variable to his book \parencite{Watanabe2009}, without a theorem number; \cref{ch:freeenergy} states it with its hypotheses (\cref{thm:freeenergy:main}) and proves it for models whose mean is linear in a vector of coefficients (\cref{thm:freeenergy:linear}). Each step of our proof has a counterpart. The reduction to one coefficient becomes resolution of singularities \parencite{Hironaka1964a}, a black box (\cref{thm:toolkit:hironaka}): in local coordinates $K=u^{2k}$ and $\varphi(w)\dd w=b(u)\abs{u^h}\dd u$, and $h$ and $k$ determine $\lambda$ and $m$. The window lemma becomes the asymptotics of the state density near $0$. The normal variable $Z_n$ becomes an empirical process $\xi_n(u)$ with $nK_n=nu^{2k}-\sqrt n\,u^k\xi_n(u)$, whose convergence in law to a Gaussian process (\cref{thm:toolkit:empirical}) needs the relatively finite variance condition. Watanabe's review \parencite{Watanabe2024} writes the limit as $\chi(\xi)=-\log\sum\iint t^{\lambda-1}e^{-t+\sqrt t\,\xi(u)}\dd t\,\dd u^*$, summed over local coordinates, with $\dd u^*$ carried by the part of the resolved space where the coordinates that attain $\lambda$ vanish. For one coefficient that part is $c_0$ approached from either side, $\dd u^*$ has mass $2^{\lambda-m}\ell_\pm$ on the two sides, $\xi=\pm\sqrt2\,Z$ there, and $t\to t^2/2$ turns his formula into \eqref{eq:coefficient:chi} (\cref{ex:coefficient:limits}). The general proof adds Hironaka's theorem, uniform control of an empirical process, and a multiplicity that counts coordinates rather than powers of a logarithm. It does not add a rate: we know of no general counterpart of \cref{prop:coefficient:rate}.

\section{Measuring it}\label{sec:coefficient:measuring}

\begin{figure}[!t]
\begin{fullwidth}
\centering
\begin{tikzpicture}
\begin{groupplot}[group style={group size=2 by 2,horizontal sep=15mm,vertical sep=13mm},splot,
  width=80mm,height=41mm,xmode=log,xmin=60,xmax=1.7e8,ymin=0,ymax=10,
  xtick={1e2,1e4,1e6,1e8},ytick={0,2,4,6,8,10}]
\nextgroupplot[title={$y=abx$: $\lambda=\frac12$, $m=2$, $d=2$},ylabel={$\E[F_n-nS_n]$},
  legend to name=coefficientfreelegend,legend columns=4,legend style={/tikz/every even column/.append style={column sep=4mm}}]
\addplot[regularpred] table[x=n,y=prodbic] {figures/data/03-coefficient-free.dat};
\addplot[theory] table[x=n,y=prodfin] {figures/data/03-coefficient-free.dat};
\addplot[theory2] table[x=n,y=prodlim] {figures/data/03-coefficient-free.dat};
\addplot[measured,ebars] table[x=n,y=prodF,y error plus=prodE,y error minus=prodE] {figures/data/03-coefficient-free.dat};
\legend{regular (BIC): $\frac d2\log n-\frac12$,singular: finite-$n$ prediction,singular: limit,measured mean}
\nextgroupplot[title={$y=a_1a_2a_3x$: $\lambda=\frac12$, $m=3$, $d=3$}]
\addplot[regularpred] table[x=n,y=chainbic] {figures/data/03-coefficient-free.dat};
\addplot[theory] table[x=n,y=chainfin] {figures/data/03-coefficient-free.dat};
\addplot[theory2] table[x=n,y=chainlim] {figures/data/03-coefficient-free.dat};
\addplot[measured,ebars] table[x=n,y=chainF,y error plus=chainE,y error minus=chainE] {figures/data/03-coefficient-free.dat};
\nextgroupplot[title={$y=ab^2x$: $\lambda=\frac14$, $m=1$, $d=2$},ylabel={$\E[F_n-nS_n]$},xlabel={sample size $n$}]
\addplot[regularpred] table[x=n,y=cuspbic] {figures/data/03-coefficient-free.dat};
\addplot[theory] table[x=n,y=cuspfin] {figures/data/03-coefficient-free.dat};
\addplot[theory2] table[x=n,y=cusplim] {figures/data/03-coefficient-free.dat};
\addplot[measured,ebars] table[x=n,y=cuspF,y error plus=cuspE,y error minus=cuspE] {figures/data/03-coefficient-free.dat};
\nextgroupplot[title={$y=ab^3x$: $\lambda=\frac16$, $m=1$, $d=2$},xlabel={sample size $n$}]
\addplot[regularpred] table[x=n,y=cubicbic] {figures/data/03-coefficient-free.dat};
\addplot[theory] table[x=n,y=cubicfin] {figures/data/03-coefficient-free.dat};
\addplot[theory2] table[x=n,y=cubiclim] {figures/data/03-coefficient-free.dat};
\addplot[measured,ebars] table[x=n,y=cubicF,y error plus=cubicE,y error minus=cubicE] {figures/data/03-coefficient-free.dat};
\end{groupplot}
\end{tikzpicture}

\vspace{1mm}
\pgfplotslegendfromname{coefficientfreelegend}
\end{fullwidth}
\caption{The normalized free energy of four one-coefficient models at the truth $0$, with inputs and noise $\Normal(0,1)$ and the uniform prior on $[-1,1]^d$. Dots: the mean of $F_n-nS_n$ over $200$ data sets at each of $13$ sample sizes, by exact quadrature, with Student-$t$ $95\%$ intervals (seeds $10^6(j+1)+1000k+r$ for model $j$, sample size $k$, data set $r$). Solid bronze: the finite-$n$ prediction, the exact representation \eqref{eq:coefficient:exact} at $S_{xx}=n$ averaged over $Z_n$. Dotted bronze: the limit $\lambda\log n-(m-1)\log\log n+\E\chi(Z)$ of \cref{thm:coefficient:free}. Dashed slate: the mean of Schwarz's BIC, $\frac d2\log n-\frac12$, which leaves the panels at $n\approx4\times10^4$ ($d=2$) and $n\approx10^3$ ($d=3$). Script: \texttt{figures/scripts/03-coefficient-free.py}.}\label{fig:coefficient:free}
\end{figure}

\begin{figure}[!t]
\centering
\begin{tikzpicture}
\begin{axis}[splot,height=0.46\linewidth,xmode=log,xmin=60,xmax=1.7e8,ymin=-0.46,ymax=0.32,
  xtick={1e2,1e4,1e6,1e8},ytick={-0.4,-0.3,-0.2,-0.1,0,0.1,0.2,0.3},
  xlabel={sample size $n$},ylabel={$\E[R_n-\chi(Z_n)]$},legend pos=south east,legend columns=2,
  legend style={font=\scriptsize}]
\addplot[guide,forget plot] coordinates {(60,0) (1.7e8,0)};
\addplot[theory] table[x=n,y=prodDfin] {figures/data/03-coefficient-free.dat};
\addplot[theory2] table[x=n,y=prodDone] {figures/data/03-coefficient-free.dat};
\addplot[theory,forget plot] table[x=n,y=chainDfin] {figures/data/03-coefficient-free.dat};
\addplot[theory2,forget plot] table[x=n,y=chainDone] {figures/data/03-coefficient-free.dat};
\addplot[theory,forget plot] table[x=n,y=cuspDfin] {figures/data/03-coefficient-free.dat};
\addplot[theory,forget plot] table[x=n,y=cubicDfin] {figures/data/03-coefficient-free.dat};
\addplot[measured,ebars] table[x=n,y=prodD,y error plus=prodDE,y error minus=prodDE] {figures/data/03-coefficient-free.dat};
\addplot[measured2,ebars] table[x=n,y=chainD,y error plus=chainDE,y error minus=chainDE] {figures/data/03-coefficient-free.dat};
\addplot[only marks,mark=diamond*,mark size=2pt,color=singulardeep,mark options={fill=white,line width=0.6pt},ebars] table[x=n,y=cuspD,y error plus=cuspDE,y error minus=cuspDE] {figures/data/03-coefficient-free.dat};
\addplot[only marks,mark=triangle*,mark size=2pt,color=muted,mark options={fill=white,line width=0.6pt},ebars] table[x=n,y=cubicD,y error plus=cubicDE,y error minus=cubicDE] {figures/data/03-coefficient-free.dat};
\legend{finite-$n$ prediction,first order $-\frac{(m-1)\gamma}{\log n}$,$y=abx$,$y=a_1a_2a_3x$,$y=ab^2x$,$y=ab^3x$}
\end{axis}
\end{tikzpicture}
\caption{The distance to the limit law, data set by data set: the mean of $R_n-\chi(Z_n)$ over the $200$ data sets of \cref{fig:coefficient:free}, with $Z_n=S_{xy}/\sqrt{S_{xx}}$ from each data set, and Student-$t$ $95\%$ intervals. \Cref{thm:coefficient:free} says that $R_n-\chi(Z_n)\to0$ in probability. Solid bronze: the finite-$n$ prediction; dotted bronze: the first-order term $-(m-1)\gamma/\log n$ of \cref{prop:coefficient:rate}, for the models with $m\ge2$. The regular line, not drawn, stays within $0.007$ of $0$. Script: \texttt{figures/scripts/03-coefficient-free.py}.}\label{fig:coefficient:gap}
\end{figure}

The script \texttt{03-coefficient-free.py} computes $F_n-nS_n$ for the five models of \cref{tab:coefficient:running} at $13$ sample sizes from $10^2$ to $10^8$, with $200$ data sets per model and sample size; at the truth $c_0=0$ the data are pure noise. For $n\le10^5$ it simulates the data sets in full, and for larger $n$ it draws $(S_{xx},Z_n)$ from their exact joint law (\cref{lem:coefficient:exact}). It evaluates \eqref{eq:coefficient:exact} by adaptive quadrature with the kit's \texttt{coefficient\_free\_energy}, to a relative error of about $10^{-10}$;\sidenote{The routine integrates over $\hat c\pm40$ window widths, splits the interval at the singular point $c=0$, where $\rho$ is infinite but integrable, and calls adaptive Gauss--Kronrod quadrature on each piece.} on simulated data at $n\le10^4$ its values agree to $10^{-11}$ with the kit's \texttt{free\_energy\_chain} and with step~5 of \cref{comp:coefficient:product} (\texttt{03-coefficient-checks.py}). Two predictions need no simulation: the limit of \cref{thm:coefficient:free}, and the \emph{finite-$n$ prediction}, which evaluates \eqref{eq:coefficient:exact} at $S_{xx}=n$ and averages over $Z_n$ by Gauss--Hermite quadrature, keeping the slowly varying factor that the limit drops.

\Cref{fig:coefficient:free} shows three things. First, BIC is wrong by a growing amount: at $n=10^8$ it charges $17.9$ nats for $y=abx$ and $27.1$ for the chain, against measured means of $6.25$ and $4.58$. Second, the finite-$n$ prediction runs through the points: $58$ of the $65$ means lie within their $95\%$ intervals of it.\sidenote{The Student-$t$ interval undercovers here, because $F_n-nS_n$ contains the skewed term $-Z_n^2/2$: in $40{,}000$ simulated sets of $200$ draws of $-Z^2/2$ its coverage was $94.3\%$, and $7$ or more misses out of $65$ then have probability $0.075$.} Third, the slopes are biased as the $\log\log n$ term predicts: straight lines in $\log n$ fitted to the means over $10^2\le n\le10^8$ have slopes $0.500$, $0.413$, $0.324$, $0.233$ and $0.159$ for the regular line, the product model, the chain and the two cusps, against $\lambda=\frac12,\frac12,\frac12,\frac14,\frac16$; for the product model and the chain the deficit is the local slope $\lambda-(m-1)/\log n$ at $n=10^5$, $0.413$ and $0.326$, and for the cusps it comes from the corrections of order $n^{-1/4}$ and $n^{-1/3}$.

\Cref{fig:coefficient:gap} shows what \cref{fig:coefficient:free} cannot: the distance $R_n-\chi(Z_n)$, which \cref{thm:coefficient:free} says tends to $0$ in probability. Subtracting $\chi(Z_n)$ data set by data set removes the noisy part of $R_n$, and since $\E\chi(Z_n)=\E\chi(Z)$, the mean is also the gap $\E[R_n]-\E\chi(Z)$. For the cusps it falls fast, from $0.262$ to $0.0075$ and from $0.125$ to $0.0012$, at the rates $n^{-1/4}$ and $n^{-1/3}$ at which their factors $\ell$ approach their limits. For the product model and the chain it falls slowly, from $-0.108$ to $-0.030$ and from $-0.399$ to $-0.065$, and tracks the finite-$n$ prediction ($60$ of $65$ means within their intervals). The first-order term of \cref{prop:coefficient:rate} describes the product model from $n\approx10^4$ on ($-0.063$ predicted, $-0.063\pm0.010$ measured) and three quarters of the chain's gap ($-0.125$ against $-0.162\pm0.023$ at $n=10^4$); the rest is the $(\log n)^{-2}$ term. No measurement contradicts the theorem.

\section{The ledger}\label{sec:coefficient:ledger}

For a model whose likelihood depends on its parameters through one scalar coefficient, the free energy is a one-dimensional integral (\cref{lem:coefficient:exact}), and its asymptotics are read off the density of the coefficient at the truth: the exponent gives $\lambda$, the power of the logarithm $m-1$ (\cref{thm:coefficient:free}). The remainder converges in law and in mean to an explicit function of one normal variable, and is within $O_p(1/\log n)$ of it when $m\ge2$ (\cref{prop:coefficient:rate}). The measurements confirm all three statements, and show that the product model's $-\log\log n$ biases its slope over every range of $n$ that can be simulated.

\ledgercard{The normalized free energy $F_n-nS_n$ of a model whose likelihood depends on $w$ through one coefficient $c=g(w)$.}
{$\frac d2\log n+O_p(1)$ (BIC): $\log n$ for $y=abx$ and $y=ab^2x$, $\frac32\log n$ for $y=a_1a_2a_3x$.}
{$\lambda\log n-(m-1)\log\log n+O_p(1)$ when the coefficient's prior density behaves like \mbox{$\abs{c-c_0}^{2\lambda-1}(\log\frac1{\abs{c-c_0}})^{m-1}$} at the truth; the remainder tends to $\chi(Z)$ in law and in mean (\cref{thm:coefficient:free}), with a $1/\log n$ correction when $m\ge2$ (\cref{prop:coefficient:rate}).}
{$y=abx$: $\lambda=\frac12$, $m=2$, $\chi(Z)=\frac12\log\frac8\pi-\frac{Z^2}2$; $y=ab^2x$: $\lambda=\frac14$, $m=1$.}
{Exact quadrature, $200$ data sets per $n$, $10^2\le n\le10^8$: slopes $0.41$, $0.32$, $0.23$, $0.16$ against $\lambda=\frac12,\frac12,\frac14,\frac16$; $\E[R_n-\chi(Z_n)]=-0.030$ for $y=abx$ and $0.0075$ for $y=ab^2x$ at $n=10^8$.}
{One scalar coefficient; Gaussian noise of known variance; a fixed realizable truth; the exact posterior at $\beta=1$; a prior law of $c$ with the stated form at $c_0$; $\E[x^2]=1$ (and $\E[x^4]<\infty$ for the rate).}

\section{Where the theory stops}\label{sec:coefficient:stops}

\paragraph{What was assumed.} Four things, and a trained network violates each. First, the likelihood depends on the parameters through one scalar: exact for $y=g(w)x$, false for the $\tanh$ neuron, whose mean is $abx$ only to first order, for the quartic model $y=ax_1+b^2x_2$, which has two coefficients, and for every network with two hidden units. (The ReLU neuron is a borderline case: $a\relu(bx)$ is $ab\relu(x)$ for $b>0$ and $-ab\relu(-x)$ for $b<0$, two one-coefficient models glued along $b=0$.) Second, Gaussian noise of known variance, which makes the energy exactly quadratic in $c$ and $Z_n$ exactly normal; an unknown variance adds $\frac12\log n$ (\cref{ex:coefficient:variance}). Third, a fixed truth: the limit holds for each $c_0$ but not uniformly in $c_0$, and a truth at distance $n^{-1/2}$ from the singular point behaves like the singular point (\cref{ex:coefficient:near}; \cref{ch:phases}). Fourth, the exact posterior at temperature $1$ under a prior whose coefficient law has the stated form at $c_0$. The prior is part of the answer: a Gaussian prior changes the constant (\cref{ex:coefficient:gaussian}), and the prior $\varphi(a,b)=\frac12\abs a$, which vanishes on the arm $a=0$ of the cross, gives $c=ab$ the density $1-\abs c$, without a logarithm, so that $m=1$.

\paragraph{Where it breaks.} The bounded term is not small, and it does not settle down at the sample sizes one can afford: in \cref{fig:coefficient:gap} the product model's remainder is still $0.030$ from its limit at $n=10^8$, and the chain's $0.065$. \Textcite{DrtonPlummer2017}, whose singular BIC uses $\lambda$ and $m$, remark that the bounded remainder is generally very difficult to estimate and that they knew of no higher-order approximations in statistically relevant singular settings. The lower-order terms matter most when learning coefficients are similar. Fitted to the same data at the truth $0$, the product model beats the regular line $y=ax$ by the log \term{Bayes factor}\index{Bayes factor} $\log(\frac12\log S_{xx}-h(Z_n))$ (\cref{ex:coefficient:bayesfactor}): on average $1.58$ nats at $n=10^4$, $2.25$ at $n=10^8$, and $3$ only at $n\approx10^{17}$. Inside one model whose prior mixes the law of the cusp's coefficient with that of the chain of length five, the posterior odds favor the chain's region, whose learning coefficient is twice as large, on average until $n\approx10^{13}$ (\cref{ex:coefficient:regions}). This is the situation of the loss--complexity trade-off that \textcite{CarrollEtAl2025} use to read the development of small transformers, in a preprint that assumes the lower-order terms of competing local free energies to cancel, noting that otherwise their difference enters the posterior odds, especially at small $n$. \Textcite{Skalse2023} argues in a blog post that the limit $n\to\infty$ is not where the generalization of deep networks needs explaining; our numbers make a narrower point: even for $y=abx$ the terms below $\lambda\log n$ still move at every practical $n$.

\paragraph{Open problems.} For one-coefficient models the free energy is completely described, exactly by \cref{lem:coefficient:exact}, in the limit by \cref{thm:coefficient:free}, and to first order in $1/\log n$ by \cref{prop:coefficient:rate}. The open questions begin with two coefficients (\cref{ch:freeenergy}); two are precise. First, compute the coefficient of $1/\log n$ in $\E[F_n-nS_n]$ for a model with several coefficients and $m\ge2$, such as rank-one regression with $M=N\ge2$ inputs and outputs; we know of no such computation. Second, \textcite{CarrollEtAl2025} leave to future work the theory of the lower-order terms of local free energies and of their role in the posterior's choice between regions; for most networks even the multiplicity of a local free energy is unknown, and an estimator of $m$ and of the constant from posterior samples, with error bars, would be a first step (\cref{ex:coefficient:regions}).

\begin{exercises}
\item \consolidation (Products of uniforms.) (a) Show that $\E\abs{a}^{2z}=(2z+1)^{-1}$ for $a$ uniform on $[-1,1]$ and $z>-\frac12$, and deduce that the zeta function of the chain is $2^{-z}(2z+1)^{-L}$. (b) Derive $\rho_L$ of \cref{comp:coefficient:densities} by induction on $L$. (c) Show that $\mu\{\abs c<r\}=r\sum_{k<L}\frac1{k!}(\log\frac1r)^k$ for $0<r\le1$, and check \cref{prop:coefficient:zeta} for the chain.\label{ex:coefficient:products}

\item \consolidation (The limit laws.) (a) Derive the values of $\chi(0)$ and $\E\chi(Z)$ in \cref{tab:coefficient:running} for the regular line, the product model and the chain. (b) Show that $\int_0^\infty t^{2\lambda-1}e^{-t^2/2+zt}\dd t=2^{\lambda-1}\int_0^\infty\tau^{\lambda-1}e^{-\tau+\sqrt{2\tau}\,z}\dd\tau$, and check the translation of Watanabe's $\chi(\xi)$ in \cref{sec:coefficient:general}.\label{ex:coefficient:limits}

\item \consolidation (The population free energy, and Gibbs.) (a) For the product model show that $\bar F(n)=-\log\int e^{-nK}\varphi$ equals $\frac12\log n-\log(\log n+\gamma+\log2)+\frac12\log\frac8\pi$ up to an exponentially small error. \hint{$\int_0^\infty e^{-t^2/2}\log t\dd t=-\frac14\sqrt{2\pi}(\gamma+\log2)$.} (b) Show that $F_n-nS_n=\min_\nu\{\E_\nu[nK_n]+\KL(\nu\,\|\,\varphi)\}$ over probability densities $\nu$ on $W$, and deduce $\E[F_n-nS_n]\le\bar F(n)$ for every model. (c) For the product model show that $\bar F(n)-\E[F_n-nS_n]\to\frac12$.\label{ex:coefficient:population}

\item \consolidation (No convergence in probability; other scales.) Let $x\sim\Normal(0,1)$. (a) Show that $(Z_n,Z_{4n})$ converges in law to a normal pair with correlation $\frac12$, and deduce that $R_n$ does not converge in probability for the product model. \hint{$T_{4n}-T_n$ is independent of $T_n$.} (b) If $\E[x^2]=\sigma_x^2$, show that $R_n\to\chi(Z)+\lambda\log\sigma_x^2$ in law. (c) At inverse temperature $\beta$, show that $F_n(\beta)-n\beta S_n-\lambda\log n+(m-1)\log\log n\to\chi(\sqrt\beta Z)+\lambda\log\beta$ in law.\label{ex:coefficient:notinprob}

\item \byhand (The family $y=a^jb^kx$.) (a) For $1\le j<k$, not both even, show that $c=a^jb^k$ has the density $\frac{\abs c^{1/k-1}-\abs c^{1/j-1}}{2(k-j)}$, so that $\lambda=\frac1{2k}$, $m=1$ and $\ell_\pm=\frac1{2(k-j)}$; check $\chi(0)=-0.446$ for the cubic cusp. What changes when $j$ and $k$ are both even? (b) For $j=k$ odd, show that $\rho(c)=\frac1{2k^2}\abs c^{1/k-1}\log\frac1{\abs c}$, so that $\lambda=\frac1{2k}$ and $m=2$.\label{ex:coefficient:family}

\item \byhand (A Gaussian prior.) Let $a,b$ be independent $\Normal(0,1)$ under the prior. (a) Show that $c=ab$ has the density $\frac1\pi K_0(\abs c)$, with $K_0(v)=\int_0^\infty e^{-v\cosh u}\dd u$. (b) Using $K_0(v)=\log\frac2v-\gamma+o(1)$, show that $\lambda=\frac12$, $m=2$, $\ell_\pm=\frac1\pi$ and $\chi(z)=\frac12\log(2\pi)-\frac{z^2}2$, which exceeds the uniform prior's limit by $\log\frac\pi2=0.45$ nats. Explain the sign.\label{ex:coefficient:gaussian}

\item \byhand (A one-sided coefficient.) Let $y=a^2x+\varepsilon$, with $a$ uniform on $[-1,1]$ and $c_0=0$. (a) Show that $\rho(c)=\frac12c^{-1/2}$ on $(0,1]$, so that $\lambda=\frac14$, $m=1$, $\ell_+=\frac12$ and $\ell_-=0$. (b) Show that $\chi(z)=-\log(\frac12\int_0^\infty e^{-t^2/2+zt}t^{-1/2}\dd t)$, with $\chi(0)=-0.075$ and $\E\chi(Z)=-0.231$. (c) Why is the cusp's $\chi(0)$ smaller by exactly $\log2$, and why does $\chi(z)$ grow as $z\to-\infty$?\label{ex:coefficient:onesided}

\item \variation\hard (Unknown noise variance.) Add a noise scale to the model, $y=g(w)x+\sigma\varepsilon$, with a prior density for $\sigma$ that is continuous and positive at $\sigma_0=1$, an interior point of its compact support. Show that $F_n-nS_n=(\lambda+\frac12)\log n-(m-1)\log\log n+O_p(1)$. \hint{integrate $\sigma$ out by Laplace's method (\cref{thm:toolkit:laplace}): up to bounded terms and terms free of $c$, this gives $\frac n2\log Q(c)+\frac12\log n$ with $Q(c)=\sum_i(y_i-cx_i)^2$, and $\frac n2\log\frac{Q(c)}{Q(\hat c)}$ is nearly the window of \cref{lem:coefficient:exact}. For $\sigma$ uniform on $[\frac12,2]$ and $\sum_i\varepsilon_i^2=n$ the extra cost is $\frac12\log\frac n\pi+\log\frac32+o(1)$.}\label{ex:coefficient:variance}

\item \variation (A truth near the singular point.) For the product model with a truth $\abs{c_0}<\frac12$: (a) show that, up to an exponentially small error, $F_n-nS_n=\frac12\log S_{xx}-\log\bigl(\frac12\log S_{xx}-h(\sqrt{S_{xx}}\,c_0+Z_n)\bigr)-\frac{Z_n^2}2+\frac12\log\frac2\pi$. (b) Deduce that $c_0=\tau n^{-1/2}$ with $\tau$ fixed gives the expansion of $c_0=0$, while a fixed $c_0\ne0$ gives $\frac12\log n-\log\log\frac1{\abs{c_0}}+\frac12\log\frac2\pi-\frac{Z_n^2}2+o_p(1)$, with $\lambda=\frac12$ and $m=1$. (c) For $c_0=0.1$ the slope of $\bar F(n)$ against $\log n$ is $0.45$, $0.53$, $0.50$ at $n=10^2,10^3,10^4$; explain the crossover near $n\approx c_0^{-2}$, and show that $\bar F(n)$ exceeds its value at $c_0=0$ ($0.82$ more at $n=10^4$) by an amount growing like $\log\log n$.\label{ex:coefficient:near}

\item \measurement (Fitting $\lambda$ and $m$.) Fit $A\log n-B\log\log n+C$ to the means computed by \texttt{03-coefficient-free.py}, weighted by their intervals: for the product model $A=0.49\pm0.01$ and $B=0.82\pm0.15$, for the chain $B=1.76\pm0.12$ (standard errors), against $m-1=1$ and $2$. With $m$ fixed, the slope of the means plus $(m-1)\log\log n$ is $0.509$ and $0.515$. Explain why $\lambda$ is well determined and $m$ is not, and which term of \cref{prop:coefficient:rate} biases $B$ downward.\label{ex:coefficient:slopes}

\item \measurement (A Bayes factor between equal learning coefficients.) Fit $y=ax$ and $y=abx$ to the same data at the truth $0$. (a) Show that $F_n(y{=}ax)-F_n(y{=}abx)=\log(\frac12\log S_{xx}-h(Z_n))$ up to an exponentially small error. (b) Compute its mean: $1.58$ nats at $n=10^4$, $2.25$ at $n=10^8$, and $3$ only at $n\approx10^{17}$. (c) What does this say about choosing between models that differ only in their multiplicity?\label{ex:coefficient:bayesfactor}

\item \research (Lower-order terms in the posterior's choice of region.) Let the prior be the equal mixture of two disjoint parameter spaces, one with $c=ab^2$ (the cusp) and one with $c=a_1\cdots a_5$ (the chain, $\lambda=\frac12$, $m=5$), each uniform. (a) Show that the posterior odds of the chain's region are $\exp(F_A-F_B)$, with $F_A$ and $F_B$ the free energies of the cusp and of the chain, and that the cusp's region wins as $n\to\infty$. (b) Compute with \eqref{eq:coefficient:exact} that $\E[F_A-F_B]$ is $0.89$, $0.19$ and $-0.37$ at $n=10^4$, $10^{12}$ and $10^{14}$. (c) Open: estimate the multiplicity and the constant of a network's local free energies from posterior samples, with error bars, and decide when they change the posterior's choice \parencite{CarrollEtAl2025}.\label{ex:coefficient:regions}
\end{exercises}

\notesandrefs
The free energy formula with the learning coefficient in place of $\frac d2$ is Watanabe's: his review \parencite{Watanabe2024} traces it to \textcite{Watanabe1999}; \textcite{Watanabe2001a} proved that the averaged free energy is $\lambda\log n-(m-1)\log\log n$ plus a constant; the book \parencite{Watanabe2009} gives the version as a random variable, which \textcite{Watanabe2013} restates as Theorem~2; and \textcite{Watanabe2018} develops its consequences. The one-coefficient theory of this chapter is a special case, with the book's own proofs. The product model is the smallest reduced-rank regression \parencite{AoyagiWatanabe2005}, and the chain the narrowest deep linear network \parencite{Aoyagi2024}.

Densities of products of independent variables are classical: \textcite{Craig1936} studied the product of two normal variables, and \textcite{Epstein1948} used the Mellin transform, which does for products what the Fourier transform does for sums. Up to the factor $2^{-z}$ and a change of variable, the zeta function of \cref{prop:coefficient:zeta} is the Mellin transform of $\rho$, which is why the chain's zeta function is the $L$th power of a uniform variable's (\cref{ex:coefficient:products}). The window heuristic is Bayesian Occam's razor: \textcite[ch.~28]{MacKay2003} writes the evidence as the best-fit likelihood times an Occam factor, the ratio of the posterior to the prior accessible volume. \Textcite{DrtonPlummer2017} turned the free energy formula into a model-selection criterion that uses both $\lambda$ and $m$; \textcite{DrtonEtAl2009} and \textcite{Lin2011} treat Bayesian integrals of algebraic models, and \textcite{WeiEtAl2023} argue for deep-learning readers that singularity is a feature of Bayesian learning, not a defect.

\paragraph{Translating the literature.} \Textcite{Watanabe2013} writes $F$ for our $F_n$, the Bayes free energy, and expands it around $nL_n(w_0)$, our $nS_n$ in the realizable case; his $R_n$ is ours. \Textcite{Watanabe2001a} calls the free energy the stochastic complexity, a name the 2009 book keeps beside the evidence (its \S6.2 is ``Evidence and stochastic complexity''). The review \parencite{Watanabe2024} writes $\theta$, $\Theta$, $\pi(\theta)$ for our $w$, $W$, $\varphi(w)$; its $t$, $\dd u^*$ and $\xi$ translate as in \cref{sec:coefficient:general}. \Textcite{DrtonPlummer2017} write $\lambda_i(\pi_0)$ and $m_i(\pi_0)$ for our $\lambda$ and $m$ of a model $\cM_i$ at the truth $\pi_0$; their proxy $P(Y_n\mid\hat\pi_i,\cM_i)(\log n)^{m_i(\pi_0)-1}/n^{\lambda_i(\pi_0)}$ for the marginal likelihood is our expansion of $e^{-F_n}$ with the remainder dropped and $e^{-nS_n}$ replaced by the maximized likelihood.

\printchapterbib
\end{refsection}
